\subsection{Simulation establishes the equivalence of HRIC turnover and PERMANOVA group separation}

We used zero-inflated microbial-community simulations to address two questions.
We first asked whether group-attributable HRIC turnover is the same geometric
quantity that drives multivariate group separation. We then evaluated the
power and type I error of the resulting PERMANOVA test under abundance,
prevalence and joint perturbations. SparseDOSSA 2 was fitted to the throat
microbiome data distributed with the \texttt{GUniFrac}
package~\cite{charlson2010disordered,chen2012associating,
ma2021sparsedossa}. After filtering, the simulation template contained 616
taxa.

The HRIC quantities occur at distinct levels. For a community table
\(X=(\pi_1,\ldots,\pi_n)^\mathsf{T}\), cohort-level beta diversity is
\[
\beta_{\mathrm{SH}}(X)
  = \gamma_{\mathrm{SH}}(X)
  - \frac{1}{n}\sum_{i=1}^{n}\alpha_{\mathrm{SH}}(\pi_i).
\]
This is one beta-diversity value for the complete cohort, rather than a
sample-to-sample dissimilarity. For a grouping variable \(g\),
\texttt{HRIC::SHdelta} instead isolates the turnover attributable to group
separation. Writing \(Z_i=\operatorname{HRIC}(\pi_i)\), it is
\[
\delta_{\mathrm{SH}}(X,g)
  = \sum_{h=1}^{G}\frac{n_h}{n}
    \left\lVert\bar Z_h-\bar Z\right\rVert_2^2
  = \frac{SS_{\mathrm{between}}}{n}.
\]
Its normalized form,
\[
\delta_{\mathrm{SH}}^*
  = \frac{\delta_{\mathrm{SH}}}{A_p^2},
\qquad
A_p=\arcsin\sqrt{1-\frac{1}{p}},
\]
therefore satisfies the exact identity
\[
SS_{\mathrm{between}}
  = nA_p^2\delta_{\mathrm{SH}}^*.
\]

This identity connects ecological turnover directly to standard multivariate
inference. For two samples, the HRIC distance is
\[
d_{\mathrm{HRIC}}(\pi_i,\pi_j)
  = \left\lVert
      \operatorname{HRIC}(\pi_i)
      - \operatorname{HRIC}(\pi_j)
    \right\rVert_2 .
\]
PERMANOVA applied to this Euclidean distance partitions the same
\(SS_{\mathrm{between}}\) that defines \(\delta_{\mathrm{SH}}\). A separate
hypothesis test for turnover is consequently unnecessary:
\texttt{HRIC::SHdelta} supplies the group-attributable turnover effect size,
whereas PERMANOVA supplies its permutation test relative to residual
within-group dispersion.

We verified this correspondence in simulations with \(k=12\)
phylogenetically coherent taxon clusters. Each cluster was perturbed in turn,
and quantities were averaged across clusters within each replicate. With
\(n=100\) samples and \(p=616\) taxa, \(nA_p^2=234.2\). The largest absolute
numerical error in
\(SS_{\mathrm{between}}=nA_p^2\delta_{\mathrm{SH}}^*\) was
\(6.1\times10^{-14}\) (Fig.~\ref{fig:simulation-hric}b). The PERMANOVA
pseudo-\(F\) statistic additionally scales this separation by within-group
dispersion:
\[
F_{\mathrm{HRIC}}
  =
  \frac{SS_{\mathrm{between}}/(G-1)}
       {SS_{\mathrm{within}}/(n-G)}.
\]
Normalized turnover was therefore strongly related, but not mathematically
identical, to pseudo-\(F\) over the simulated range
(\(R^2=0.998,0.999\) and \(0.993\) for abundance, joint and prevalence
perturbations, respectively; Fig.~\ref{fig:simulation-hric}a). Thus, the
equivalence is exact for turnover and the PERMANOVA numerator. The denominator
then standardizes that turnover by residual community heterogeneity.

Having established this equivalence, we evaluated the power of PERMANOVA
without introducing an additional test. The reference phylogeny was
partitioned into \(k=2,4,6\) or 8 coherent and approximately size-balanced
clusters, and one cluster at a time was perturbed in abundance, prevalence or
both. Each dataset contained 100 samples independently assigned to one of two
conditions with probability 0.5. HRIC distance was compared with
Bray--Curtis, ordinary Euclidean, presence--absence Jaccard, Aitchison and
\(\operatorname{ILR}\)--Euclidean distances (Fig.~\ref{fig:simulation-hric}c).
Under abundance-only perturbations, HRIC distance closely tracked
Bray--Curtis and reached 80\% empirical power at effect sizes of
\(0.75,0.75,0.75\) and \(0.80\) for \(k=2,4,6\) and 8, respectively.
The corresponding thresholds for Bray--Curtis were
\(0.70,0.75,0.75\) and \(0.70\), whereas none of the displayed Aitchison,
\(\operatorname{ILR}\)--Euclidean or Jaccard variants reached 80\% power.
Under joint perturbations, HRIC distance reached 80\% power at effect sizes of
\(0.50\)--\(0.55\), with several log-ratio and occurrence-sensitive
distances performing similarly or detecting smaller effects. The ordering
reversed under prevalence-only perturbations. At an effect size of \(0.50\),
power for HRIC distance was \(0.07\)--\(0.11\), compared with
\(0.48\)--\(1.00\) for Jaccard and \(0.39\)--\(0.79\) for the Aitchison
distance with a pseudocount of 0.5. HRIC distance was therefore most
responsive when turnover involved redistribution of compositional mass, but
was not designed as a pure occurrence test.

At effect size zero, the mean type I error of the HRIC-distance PERMANOVA was
0.0533 across the 12 combinations of signal mode and \(k\), with a range of
0.0200--0.0900 (Table~\ref{tab:simulation-type1}). Across all displayed
methods, overall mean type I error ranged from 0.0367 to 0.0600. The wider
setting-specific range of 0.0100--0.1200 occurred in individual
100-replicate cells, whereas the method-level means showed no systematic
inflation.

\begin{figure*}[t]
    \centering
    \includegraphics[
        width=\textwidth,
        height=0.82\textheight,
        keepaspectratio
    ]{figs/figure_integrated_simulation_hric_turnover.pdf}
    \caption{\textbf{Geometric equivalence of HRIC turnover and PERMANOVA
    separation, followed by signal-dependent power.}
    \textbf{a}, Relationship between normalized between-group turnover
    \(\delta_{\mathrm{SH}}^*
    =\texttt{HRIC::SHdelta}(X,g)/A_p^2\) and the trace-based HRIC
    pseudo-\(F\) statistic in the \(k=12\) simulation. Each signal mode and
    effect size comprised 125 independent replicates. Metrics were averaged
    across the 12 possible signal clusters within each replicate. Points are
    replication-level estimates; lines and shaded bands are mode-specific
    ordinary least-squares fits and 95\% confidence intervals, with \(R^2\)
    shown in each facet. \textbf{b}, Exact relationship between normalized
    between-group turnover and \(SS_{\mathrm{between}}\). Points are
    replication-level estimates and the black dashed line is the theoretical
    identity. For \(n=100\) and \(p=616\), \(nA_p^2=234.2\); the largest
    absolute numerical error in
    \(SS_{\mathrm{between}}=nA_p^2\delta_{\mathrm{SH}}^*\) was
    \(6.1\times10^{-14}\). \textbf{c}, Power of PERMANOVA based on HRIC
    distance and comparator dissimilarities under perturbations to taxon
    abundance, abundance and prevalence jointly, or prevalence alone. Taxa
    were divided into \(k=2,4,6\) or 8 phylogenetically coherent clusters.
    For each \(k\), the displayed signal cluster contained the largest number
    of taxa (\(308,154,103\) and 77 taxa, respectively); ties were resolved by
    total abundance in the fitted template. Each setting comprised 100
    independently simulated datasets of 100 samples. Curves are binomial
    generalized additive model fits to empirical rejection proportions, and
    horizontal grey lines mark 80\% power. HRIC distance is
    \(d_{\mathrm{HRIC}}(\pi_i,\pi_j)
    =\lVert\operatorname{HRIC}(\pi_i)-
    \operatorname{HRIC}(\pi_j)\rVert_2\). Parenthetical values for Aitchison
    and \(\operatorname{ILR}\)--Euclidean distances denote the pseudocount
    added before log-ratio transformation.}
    \label{fig:simulation-hric}
\end{figure*}

\begin{table*}[t]
    \centering
    \caption{\textbf{Empirical type I error of distance-based PERMANOVA in
    null simulations.}
    Values under each nominal signal mode are mean rejection proportions
    across \(k=2,4,6\) and 8 at effect size zero. The final column reports the
    overall mean and range across all 12 combinations of mode and \(k\).
    Each combination contained 100 independent simulated datasets, and tests
    used a nominal level of 0.05. Parenthetical values in method names denote
    the pseudocount added before log-ratio transformation.}
    \label{tab:simulation-type1}
    \begin{tabular}{lcccc}
        \toprule
        Method & Abundance & Joint & Prevalence & Overall mean (range) \\
        \midrule
        HRIC distance & 0.0525 & 0.0500 & 0.0575 &
        0.0533 (0.0200--0.0900) \\
        Bray--Curtis & 0.0600 & 0.0450 & 0.0500 &
        0.0517 (0.0200--0.1000) \\
        Euclidean & 0.0450 & 0.0300 & 0.0350 &
        0.0367 (0.0100--0.0700) \\
        Jaccard & 0.0375 & 0.0425 & 0.0700 &
        0.0500 (0.0100--0.0800) \\
        Aitchison (0.01) & 0.0425 & 0.0550 & 0.0675 &
        0.0550 (0.0200--0.1000) \\
        Aitchison (0.1) & 0.0500 & 0.0600 & 0.0625 &
        0.0575 (0.0100--0.1000) \\
        Aitchison (0.5) & 0.0500 & 0.0600 & 0.0600 &
        0.0567 (0.0200--0.1000) \\
        Aitchison (1) & 0.0525 & 0.0675 & 0.0500 &
        0.0567 (0.0200--0.1000) \\
        \(\operatorname{ILR}\)--Euclidean (0.01) &
        0.0475 & 0.0600 & 0.0625 & 0.0567 (0.0200--0.1000) \\
        \(\operatorname{ILR}\)--Euclidean (0.1) &
        0.0525 & 0.0550 & 0.0650 & 0.0575 (0.0200--0.1100) \\
        \(\operatorname{ILR}\)--Euclidean (0.5) &
        0.0550 & 0.0625 & 0.0625 & 0.0600 (0.0200--0.1200) \\
        \(\operatorname{ILR}\)--Euclidean (1) &
        0.0500 & 0.0600 & 0.0500 & 0.0533 (0.0200--0.1000) \\
        \bottomrule
    \end{tabular}
\end{table*}

% Methods prose to incorporate into the manuscript's continuous Methods section.

We fitted SparseDOSSA 2 (v.0.99.2) to the throat 16S rRNA gene OTU table and
phylogeny distributed with \texttt{GUniFrac}~\cite{charlson2010disordered,
chen2012associating,ma2021sparsedossa}. Taxa occurring in fewer than two
reference samples were removed, leaving 616 taxa. Patristic distances were
calculated from the pruned reference tree, and balanced polar-ordination
clustering was used to construct phylogenetically coherent groups of taxa.
Binary condition labels were drawn independently from a Bernoulli
distribution with probability 0.5. A positive SparseDOSSA 2 effect shifted
the fitted log-scale abundance mean, the log odds of taxon presence, or both
for every taxon in the selected cluster.

For the turnover--PERMANOVA correspondence analysis, we set \(k=12\) and used
each cluster in turn as the perturbed set. We generated 125 independent
replicates for each of 21 effect sizes and each signal mode. Within each
replicate, quantities were averaged across the 12 signal clusters. In current
package notation, normalized between-group turnover was
\(\delta_{\mathrm{SH}}^*
=\texttt{HRIC::SHdelta}(X,g)/A_p^2\). For
\(Z_i=\operatorname{HRIC}(\pi_i)\), we calculated
\[
SS_{\mathrm{between}}
  = \sum_{h=1}^{G}n_h
    \left\lVert\bar Z_h-\bar Z\right\rVert_2^2
\]
and
\[
SS_{\mathrm{within}}
  = \sum_{h=1}^{G}\sum_{i:g_i=h}
    \left\lVert Z_i-\bar Z_h\right\rVert_2^2 .
\]
The trace-based pseudo-\(F\) statistic was
\((SS_{\mathrm{between}}/(G-1))/(SS_{\mathrm{within}}/(n-G))\).
Panel a shows mode-specific ordinary least-squares fits with 95\% confidence
intervals. Panel b compares replication-level estimates with the exact
identity
\(SS_{\mathrm{between}}=nA_p^2\delta_{\mathrm{SH}}^*\).

For the power analysis, \(k\in\{2,4,6,8\}\), and every cluster was used in
turn as the perturbed set. Figure~\ref{fig:simulation-hric}c shows, for each
\(k\), the cluster containing the largest number of taxa; size ties were
resolved using fitted abundance rank. For each signal cluster, signal mode
and effect size, we generated 100 independent datasets containing 100
samples. The effect parameter ranged from 0 to 1 in increments of 0.05.
HRIC coordinates were calculated directly with \texttt{HRIC::HRIC} (v.0.1.0)
and their Euclidean distances were used without zero replacement.
Community-level association was tested with \texttt{vegan::adonis2}
(v.2.7-5) using 999 permutations~\cite{anderson2001new}. Comparator
distances were Bray--Curtis, ordinary Euclidean and binary Jaccard, together
with Aitchison and \(\operatorname{ILR}\)--Euclidean distances after adding
pseudocounts of 0.01, 0.1, 0.5 or 1. Weighted, unweighted and generalized
UniFrac variants were also evaluated but are not displayed in the main
figure. Power was the proportion of replicate \(P\) values below 0.05 for
non-zero effects. For visualization, each empirical power series was fitted
with a binomial generalized additive model using a shrinkage cubic regression
spline, and the displayed fit was constrained to be non-decreasing over
increasing effect size.

Type I error was the proportion of permutation \(P\) values below 0.05 at
effect size zero. To match panel c, the largest taxon cluster was selected
within every combination of \(k\) and nominal signal mode, with fitted
abundance rank used to break size ties. Every setting contained 100
independent null datasets. Table~\ref{tab:simulation-type1} reports the mean
type I error across the four values of \(k\) within each mode, together with
the overall mean and range across the 12 mode-by-\(k\) settings.
